\begin{table}[htbp]
\centering
\caption{Comprehensive Ablation Analysis: Impact of Stability Gating Thresholds and Feature Classes on Downstream Model Calibration.}
\label{tab:ablation}
\small
\setlength{\tabcolsep}{4.8pt}
\begin{tabular}{lcccccc}
\toprule
\textbf{Ablation Condition} & $\mathbf{N_{\text{feats}}}$ & $\mathbf{AUC_{\text{base}}}$ & $\mathbf{Brier_{\text{base}}}$ & $\mathbf{Brier_{-2\text{vox}}}$ & $\mathbf{Brier_{\text{nnU-Net}}}$ & $\mathbf{\Delta Brier_{\text{rel}}}$ (\%) \\
\midrule
\multicolumn{7}{l}{\textit{Part A: Stability Gating Cutoff ($\tau = ICC_{\text{overall}}$)}} \\
$\tau \ge 0.00$ (Full Feature Pool) & 107 & \textbf{0.986} & \textbf{0.065} & 0.371 & 0.142 & +470.7\% \\
$\tau \ge 0.60$ & 44 & 0.944 & 0.102 & 0.305 & \textbf{0.140} & +198.4\% \\
$\tau \ge 0.70$ (Tier 2 Cutoff) & 36 & 0.850 & 0.161 & \textbf{0.173} & 0.199 & +23.5\% \\
$\tau \ge 0.75$ & 35 & 0.850 & 0.161 & \textbf{0.173} & 0.199 & +23.5\% \\
$\tau \ge 0.80$ & 28 & 0.606 & 0.234 & 0.242 & 0.248 & \textbf{+6.0\%} \\
$\tau \ge 0.85$ (Adopted Tier 1) & 23 & 0.624 & 0.232 & 0.243 & 0.247 & +6.7\% \\
$\tau \ge 0.90$ & 17 & 0.622 & 0.232 & 0.243 & 0.247 & +6.7\% \\
\midrule
\multicolumn{7}{l}{\textit{Part B: Leave-One-Class-Out Feature Ablation}} \\
All 107 Features (Baseline Pool) & 107 & 0.986 & 0.065 & 0.371 & 0.142 & +470.7\% \\
w/o Shape Features & 93 & 0.978 & 0.075 & \textbf{0.335} & 0.142 & \textbf{+348.8\%} \\
w/o First-Order Features & 89 & 0.982 & 0.070 & 0.381 & 0.144 & +447.7\% \\
w/o GLCM Features & 83 & 0.986 & 0.064 & 0.373 & 0.142 & +482.6\% \\
w/o GLDM Features & 93 & 0.986 & 0.064 & 0.372 & 0.141 & +480.1\% \\
w/o GLRLM Features & 91 & 0.986 & 0.065 & 0.371 & 0.142 & +472.5\% \\
w/o GLSZM Features & 91 & 0.986 & 0.065 & 0.372 & 0.142 & +475.9\% \\
w/o NGTDM Features & 102 & \textbf{0.988} & \textbf{0.057} & 0.373 & \textbf{0.133} & +556.1\% \\
\bottomrule
\end{tabular}
\begin{flushleft}
\footnotesize \textit{Note:} Models were trained on baseline ground truth expert segmentations using 5-fold cross-validation. Relative calibration drift is defined as $\Delta Brier_{\text{rel}} = \frac{\max(|Brier_{-2\text{vox}} - Brier_{\text{base}}|, |Brier_{\text{nnU-Net}} - Brier_{\text{base}}|)}{Brier_{\text{base}}} \times 100\%$. The inflection point at $\tau \ge 0.85$ provides mathematical justification for the Tier 1 cutoff, reducing calibration drift from $+470.7\%$ to $+6.7\%$. In Part B, omitting shape features yields the largest single reduction in calibration drift (from $+470.7\%$ to $+348.8\%$), demonstrating that geometry metrics constitute the primary driver of calibration failure. Bold values denote optimal metrics per column within each ablation panel.
\end{flushleft}
\end{table}
